\documentclass[12pt, a4paper]{article}
% --- 1. Cấu hình Tiếng Việt và Font chữ chuẩn ---
\usepackage[utf8]{inputenc}
\usepackage[T5]{fontenc}
\usepackage[vietnamese]{babel}
\usepackage{mathptmx} % Font Times New Roman đồng bộ cả chữ và công thức toán, không gây xung đột gói

\makeatletter
\renewcommand{\normalsize}{\@setfontsize\normalsize{13pt}{16pt}}
\makeatother
\normalsize

% --- 2. Gói Toán học & Ký hiệu bổ trợ ---
\usepackage{amsmath, amssymb, amsfonts, mathrsfs, amsthm}
\usepackage{graphicx}
\usepackage{fontawesome}
\usepackage{booktabs, tabularx, array, multirow}
\usepackage{enumitem}

% --- 3. Đồ họa TikZ, Bảng biến thiên & Hình không gian ---
\usepackage{tikz}
\usetikzlibrary{calc, intersections, angles, quotes, patterns, arrows.meta, 3d}
\usepackage{pgfplots}
\pgfplotsset{compat=1.18}
\usepackage{tkz-euclide}
\usepackage{tkz-tab}

\renewcommand{\vec}[1]{\overrightarrow{#1}}

% --- 4. Hộp sư phạm (Tcolorbox) ---
\usepackage[many]{tcolorbox}
\tcbuselibrary{skins, breakable}

% --- 5. Căn lề chuẩn văn bản giáo án ---
\usepackage{geometry}
\geometry{
    a4paper,
    left=25mm,
    right=15mm,
    top=20mm,
    bottom=20mm
}

% --- 6. Header / Footer chuẩn hóa ---
\usepackage{fancyhdr}
\usepackage{lastpage}
\pagestyle{fancy}
\fancyhf{}

\fancyhead[L]{\small \textit{Kế hoạch bài dạy Toán 11}}
\fancyhead[R]{\textbf{Trang \thepage/\pageref{LastPage}}}
\fancyfoot[L]{\small \textit{Tổ Toán - Tin}}
\fancyfoot[R]{\small \textit{Trường THCS\&THPT Hồng Vân}}
\renewcommand{\headrulewidth}{0.4pt}
\renewcommand{\footrulewidth}{0.4pt}

% --- 7. Giãn dòng & Giãn đoạn theo chuẩn Nghị định 30/2020/NĐ-CP ---
\usepackage{setspace}
\setstretch{1.15}
\setlength{\parindent}{1.0cm}
\setlength{\parskip}{5pt}

% Định nghĩa hộp sư phạm chuyên dụng
\newtcolorbox{pedabox}[2][]{
    enhanced,
    breakable,
    colback=blue!3!white,
    colframe=blue!65!black,
    fonttitle=\bfseries\small,
    title={#2},
    arc=2mm,
    boxrule=0.6pt,
    left=3mm, right=3mm, top=2mm, bottom=2mm,
    #1
}

\newtcolorbox{aibox}[2][]{
    enhanced,
    breakable,
    colback=teal!4!white,
    colframe=teal!70!black,
    fonttitle=\bfseries\small,
    title={#2},
    arc=2mm,
    boxrule=0.6pt,
    left=3mm, right=3mm, top=2mm, bottom=2mm,
    #1
}

\newtcolorbox{scaffoldbox}[2][]{
    enhanced,
    breakable,
    colback=orange!4!white,
    colframe=orange!80!black,
    fonttitle=\bfseries\small,
    title={#2},
    arc=2mm,
    boxrule=0.6pt,
    left=3mm, right=3mm, top=2mm, bottom=2mm,
    #1
}

\begin{document}

\begin{center}
    {\fontsize{14pt}{17pt}\selectfont \textbf{KẾ HOẠCH BÀI DẠY MÔN TOÁN LỚP 11}}\\[6pt]
    {\fontsize{16pt}{19pt}\selectfont \textbf{CHUYÊN ĐỀ 1 - BÀI 1. PHÉP BIẾN HÌNH}}\\[4pt]
    {\fontsize{12pt}{15pt}\selectfont \textbf{Môn học: Chuyên đề học tập Toán 11} \ \textbar \ \textbf{Thời lượng: 2 tiết (Tiết CĐ1, CĐ2 theo PPCT)}}
\end{center}

\vspace{5pt}

\section*{I. MỤC TIÊU}

\subsection*{1. Kiến thức, Năng lực}

\begin{itemize}[leftmargin=1.2cm]
    \item \textbf{Yêu cầu cần đạt (Nguyên văn CT GDPT 2018):}
    \begin{itemize}[leftmargin=0.6cm]
        \item Nhận biết được khái niệm phép biến hình, phép dời hình.
        \item Vận dụng được các phép dời hình trong đồ họa và trong một số vấn đề thực tiễn (ví dụ: tạo các hoa văn, hình khối,...).
    \end{itemize}

    \item \textbf{Năng lực Toán học đặc thù:}
    \begin{itemize}[leftmargin=0.6cm]
        \item \textit{Năng lực tư duy và lập luận toán học:} So sánh được sự khác biệt giữa quy tắc là phép biến hình (mỗi điểm có duy nhất một ảnh) và quy tắc không phải phép biến hình; phân tích và lập luận tính chất bảo toàn khoảng cách của phép dời hình; chứng minh phép dời hình bảo toàn tính thẳng hàng, góc, độ dài đoạn thẳng và diện tích tam giác.
        \item \textit{Năng lực mô hình hóa toán học:} Thiết lập mô hình phép biến hình để giải thích các hiện tượng thực tế: hình ảnh qua gương phẳng, bóng nắng của cọc tiêu trên sân trường, hoa văn lặp lại trên vải thổ cẩm của đồng bào dân tộc thiểu số A Lưới.
        \item \textit{Năng lực giải quyết vấn đề toán học:} Xác định tọa độ điểm ảnh, phương trình đường thẳng ảnh, đường tròn ảnh qua các phép biến hình cơ bản trong mặt phẳng tọa độ $Oxy$.
        \item \textit{Năng lực giao tiếp toán học:} Sử dụng chính xác các thuật ngữ sư phạm: phép biến hình, điểm gốc, điểm ảnh, phép dời hình, bảo toàn khoảng cách, phép đồng nhất.
        \item \textit{Năng lực sử dụng công cụ, phương tiện học toán:} Sử dụng compa, ê-ke để vẽ hình ảnh đối xứng; khai thác phần mềm hình học động GeoGebra để trực quan hóa phép dời hình.
    \end{itemize}

    \item \textbf{Năng lực số (Theo Thông tư 02/2025/TT-BGDĐT):}
    \begin{itemize}[leftmargin=0.6cm]
        \item \textit{Mã 3.1NC1a:} Học sinh sử dụng thành thạo phần mềm GeoGebra tạo phép biến hình động, quan sát sự thay đổi vị trí của điểm ảnh $M'$ khi kéo rê điểm gốc $M$.
        \item \textit{Mã 1.1NC1a:} Học sinh tìm kiếm, chọn lọc video thực tế trên Internet về ứng dụng phép biến hình trong cơ khí chế tạo, robot công nghiệp và đồ họa 3D.
    \end{itemize}

    \item \textbf{Năng lực AI (Theo Quyết định 2422/QĐ-BGDĐT):}
    \begin{itemize}[leftmargin=0.6cm]
        \item \textit{Mã 11.C2.1:} Học sinh trình bày và giải thích được cách thuật toán trí tuệ nhân tạo (AI Computer Vision) xử lý hình ảnh, nhận diện vật thể và biến đổi tọa độ điểm ảnh trong không gian.
        \item \textit{Mã 11.A1.2:} Học sinh phân tích được tầm quan trọng của việc kiểm soát con người khi AI dựng hình đồ họa tự động (tránh hiện tượng ảo giác AI -- Hallucination hoặc sai lệch hình học).
    \end{itemize}

    \item \textbf{Năng lực chung:}
    \begin{itemize}[leftmargin=0.6cm]
        \item \textit{Tự chủ và tự học:} Tự giác nghiên cứu tài liệu chuyên đề học tập, hoàn thành các bài tập phân hóa trên phiếu học tập cá nhân.
        \item \textit{Giao tiếp và hợp tác:} Tương tác chủ động, phân công công việc cụ thể khi thảo luận nhóm và thực hành phần mềm.
    \end{itemize}
\end{itemize}

\subsection*{2. Phẩm chất}

\begin{itemize}[leftmargin=1.2cm]
    \item \textbf{Chăm chỉ:} Kiên trì vẽ hình chính xác, tích cực thực hành phần mềm toán học và tìm kiếm tư liệu số.
    \item \textbf{Trung thực:} Báo cáo chính xác kết quả thảo luận nhóm và bài tập thực hành cá nhân.
    \item \textbf{Trách nhiệm:} Hỗ trợ các bạn học sinh trung bình -- yếu trong nhóm hiểu rõ quy tắc tương ứng một -- một của phép biến hình; bảo vệ và tự hào về bản sắc hoa văn truyền thống dân tộc.
\end{itemize}

\section*{II. THIẾT BỊ DẠY HỌC VÀ HỌC LIỆU}

\begin{itemize}[leftmargin=1.2cm]
    \item \textbf{Giáo viên:}
    \begin{itemize}[leftmargin=0.6cm]
        \item Kế hoạch bài dạy chuyên đề, bài trình chiếu đa phương tiện (PowerPoint) tích hợp video ứng dụng thực tế.
        \item File mô phỏng GeoGebra trực quan hóa phép biến hình: điểm $M \mapsto M'$ và bảo toàn khoảng cách $M'N' = MN$.
        \item Hình ảnh chụp hoa văn thổ cẩm A Lưới, hình ảnh bóng nắng cọc tiêu sân trường, gương phẳng.
        \item Hệ thống Phiếu học tập phân tầng (Worksheet CĐ 1, Worksheet CĐ 2) và Rubric đánh giá năng lực 4 mức độ.
    \end{itemize}
    \item \textbf{Học sinh:}
    \begin{itemize}[leftmargin=0.6cm]
        \item Sách Chuyên đề học tập Toán 11, vở ghi, giấy nháp.
        \item Thước thẳng có vạch chia, compa, ê-ke, bút màu.
        \item Điện thoại thông minh hoặc máy tính bảng/laptop cài đặt sẵn ứng dụng GeoGebra Geometry.
    \end{itemize}
\end{itemize}

\section*{III. TIẾN TRÌNH DẠY HỌC}

% =========================================================================
% TIẾT CĐ1
% =========================================================================
\subsection*{TIẾT CĐ1. KHÁI NIỆM PHÉP BIẾN HÌNH VÀ ẢNH CỦA MỘT HÌNH}
\textit{(Tiết Chuyên đề 1 theo PPCT | Tuần 1)}

\subsubsection*{1. Hoạt động 1: Khởi động (Mở đầu)}
\textbf{a) Mục tiêu:} Kích thích trí tò mò của học sinh về mối quan hệ giữa một điểm gốc và điểm ảnh tương ứng thông qua các hiện tượng thực tế (gương phẳng, bóng nắng, hình chiếu).

\textbf{b) Nội dung:} Giáo viên trình chiếu hai tình huống thực tế:
\begin{itemize}[leftmargin=0.6cm]
    \item \textbf{Tình huống 1:} Một học sinh đứng trước một tấm gương phẳng. Ứng với mỗi điểm $M$ trên cơ thể học sinh, có xác định được điểm ảnh $M'$ tương ứng trong gương không? Có thể có một điểm $M$ mà tạo ra hai điểm ảnh khác nhau trong gương phẳng không?
    \item \textbf{Tình huống 2:} Giữa trưa nắng, chiếc cọc tiêu cắm thẳng đứng trên sân trường tạo bóng trên mặt đất. Ứng với mỗi đỉnh $A$ của cọc tiêu có mấy bóng điểm $A'$ trên mặt sân?
\end{itemize}

\textbf{c) Sản phẩm:} Học sinh trả lời:
\begin{itemize}[leftmargin=0.6cm]
    \item Tình huống 1: Mỗi điểm $M$ chỉ cho \textbf{duy nhất một} điểm ảnh $M'$ trong gương phẳng. Không thể có một điểm mà cho hai ảnh khác nhau.
    \item Tình huống 2: Đỉnh $A$ của cọc chỉ cho \textbf{duy nhất một} bóng điểm $A'$ trên mặt đất tại một thời điểm xác định.
\end{itemize}

\textbf{d) Tổ chức thực hiện:}
\begin{itemize}[leftmargin=0.6cm]
    \item \textit{Chuyển giao nhiệm vụ:} Giáo viên chiếu hình ảnh, đặt câu hỏi cho 35 học sinh làm việc theo cặp đôi (trong 2 phút).
    \item \textit{Thực hiện nhiệm vụ (Giàn giáo cho HS TB-Yếu):}
    \begin{scaffoldbox}{Giàn giáo sư phạm CĐ1.1: Nhận diện quy tắc tương ứng}
        \textbf{Câu hỏi định hướng:}
        \begin{enumerate}[leftmargin=0.5cm]
            \item Ở lớp 10, khái niệm hàm số $y = f(x)$ yêu cầu mỗi giá trị $x$ tương ứng với bao nhiêu giá trị $y$? (Trả lời: Duy nhất một giá trị $y$).
            \item Vậy nếu thay số $x$ bằng điểm $M$ trong mặt phẳng và số $y$ bằng điểm $M'$, quy tắc này gọi là gì?
        \end{enumerate}
    \end{scaffoldbox}
    \item \textit{Báo cáo, thảo luận:} 2 cặp đôi trả lời, các cặp khác lắng nghe và nhận xét.
    \item \textit{Kết luận, nhận định:} Giáo viên chốt: Quy tắc đặt tương ứng mỗi điểm $M$ với duy nhất một điểm $M'$ trong mặt phẳng được gọi là \textbf{Phép biến hình}.
\end{itemize}

\subsubsection*{2. Hoạt động 2: Hình thành kiến thức}

\paragraph{Nội dung 1: Định nghĩa Phép biến hình và Phép đồng nhất}
\textbf{a) Mục tiêu:} Nhận biết định nghĩa phép biến hình trong mặt phẳng; phân biệt được quy tắc là phép biến hình và quy tắc không phải phép biến hình; hiểu khái niệm phép đồng nhất.

\textbf{b) Nội dung:}
\begin{itemize}[leftmargin=0.6cm]
    \item \textbf{Định nghĩa:} Quy tắc đặt tương ứng mỗi điểm $M$ trong mặt phẳng với một điểm xác định \textbf{duy nhất} $M'$ của mặt phẳng đó được gọi là một \textbf{phép biến hình} trong mặt phẳng.
    \item Ký hiệu phép biến hình là $f$:
    $$f(M) = M' \quad \text{hoặc} \quad f: M \mapsto M'$$
    Điểm $M'$ được gọi là \textbf{ảnh} của điểm $M$ qua phép biến hình $f$.
    \item \textbf{Phép đồng nhất:} Phép biến hình biến mỗi điểm $M$ thành chính nó được gọi là \textbf{phép đồng nhất}, ký hiệu là $e$ hoặc $I$:
    $$e(M) = M \quad (\forall M)$$
\end{itemize}

\textbf{c) Sản phẩm:} Học sinh lấy được ví dụ cụ thể về phép biến hình và phản ví dụ:
\begin{itemize}[leftmargin=0.6cm]
    \item \textit{Ví dụ 1 (Phép biến hình):} Trong mặt phẳng tọa độ $Oxy$, quy tắc $f$ biến điểm $M(x; y)$ thành điểm $M'(x + 2; y - 3)$ là một phép biến hình vì với mỗi cặp $(x; y)$, ta xác định được duy nhất một cặp tọa độ $(x + 2; y - 3)$.
    \item \textit{Ví dụ 2 (Không phải phép biến hình):} Quy tắc biến mỗi điểm $M$ thành điểm $M'$ sao cho khoảng cách $MM' = 2\text{ cm}$ \textbf{không phải} là phép biến hình, vì có vô số điểm $M'$ nằm trên đường tròn tâm $M$ bán kính $2\text{ cm}$ (không thỏa mãn tính duy nhất).
\end{itemize}

\textbf{d) Tổ chức thực hiện:}
\begin{itemize}[leftmargin=0.6cm]
    \item \textit{Chuyển giao nhiệm vụ:} Giáo viên trình bày định nghĩa, yêu cầu học sinh làm việc nhóm đôi phân tích Ví dụ 1 và Ví dụ 2.
    \item \textit{Thực hiện nhiệm vụ (Giàn giáo cho HS TB-Yếu):}
    \begin{scaffoldbox}{Giàn giáo sư phạm CĐ1.2: Dấu hiệu cốt lõi nhận biết phép biến hình}
        \textbf{Khắc sâu tiêu chí ''DUY NHẤT'':}
        \begin{itemize}[leftmargin=0.5cm]
            \item Có điểm gốc $M \implies$ bắt buộc tìm được ảnh $M'$.
            \item Điểm ảnh $M'$ phải là \textbf{duy nhất 01 điểm}. Nếu tìm được từ 2 điểm ảnh trở lên, hoặc không tìm được điểm ảnh nào $\implies$ \textbf{KHÔNG PHẢI PHÉP BIẾN HÌNH}.
        \end{itemize}
    \end{scaffoldbox}
    \item \textit{Báo cáo, thảo luận:} 2 học sinh lên bảng giải thích tại sao phép chiếu vuông góc lên đường thẳng $\Delta$ là phép biến hình, còn phép lấy điểm đối xứng qua hai đường thẳng song song cùng lúc thì không duy nhất.
    \item \textit{Kết luận, nhận định:} Giáo viên chuẩn hóa định nghĩa, lưu ý phép biến hình là sự tổng quát hóa của khái niệm hàm số từ tập hợp số sang tập hợp điểm.
\end{itemize}

\paragraph{Nội dung 2: Khái niệm ảnh của một hình qua phép biến hình}
\textbf{a) Mục tiêu:} Nhận biết định nghĩa ảnh của một hình qua phép biến hình; biết cách xác định ảnh của một đoạn thẳng, tam giác, đường tròn qua phép biến hình đơn giản; thực hành phần mềm GeoGebra.

\textbf{b) Nội dung:}
\begin{itemize}[leftmargin=0.6cm]
    \item \textbf{Định nghĩa ảnh của hình:} Giả sử $f$ là một phép biến hình và $H$ là một hình trong mặt phẳng. Tập hợp tất cả các điểm $M' = f(M)$ với mọi $M \in H$ được gọi là \textbf{ảnh của hình $H$} qua phép biến hình $f$, ký hiệu là $H' = f(H)$:
    $$H' = f(H) = \{M' = f(M) \mid M \in H\}$$
    \item Để xác định ảnh của một hình, ta xác định ảnh của các điểm đặc trưng cấu tạo nên hình đó (ví dụ: các đỉnh của đa giác, tâm và bán kính của đường tròn).
\end{itemize}

\textbf{c) Sản phẩm:} Học sinh xác định được ảnh của tam giác $ABC$ qua phép tịnh tiến hoặc phép chiếu; thực hành vẽ trên GeoGebra ra hình ảnh trực quan.

\textbf{d) Tổ chức thực hiện (Tích hợp Năng lực số 3.1NC1a):}
\begin{itemize}[leftmargin=0.6cm]
    \item \textit{Chuyển giao nhiệm vụ:} Giáo viên yêu cầu học sinh mở phần mềm GeoGebra trên điện thoại/laptop, vẽ tam giác $ABC$ và tạo điểm $M'$ dịch chuyển theo vectơ cố định $\vec{u} = (3; 2)$.
    \item \textit{Thực hiện nhiệm vụ (Năng lực số):}
    \begin{aibox}{Thực hành Năng lực số: Tạo phép biến hình động trên GeoGebra}
        \textbf{Các bước thao tác:}
        \begin{enumerate}[leftmargin=0.5cm]
            \item Chọn công cụ \texttt{Polygon} $\to$ Vẽ tam giác $ABC$.
            \item Tạo vectơ $\vec{u}$ bằng công cụ \texttt{Vector}.
            \item Dùng công cụ \texttt{Translate by Vector} $\to$ Nhấp vào tam giác $ABC$ rồi nhấp vào vectơ $\vec{u} \to$ Xuất hiện tam giác ảnh $A'B'C'$.
            \item Kéo rê các đỉnh $A, B, C$ hoặc thay đổi vectơ $\vec{u} \to$ Quan sát tam giác ảnh $A'B'C'$ chuyển động tương ứng một -- một.
        \end{enumerate}
    \end{aibox}
    \item \textit{Báo cáo, thảo luận:} 2 học sinh chiếu màn hình thao tác GeoGebra lên tivi lớp học, nhận xét về sự biến thiên đồng bộ giữa hình gốc và hình ảnh.
    \item \textit{Kết luận, nhận định:} Giáo viên nhận xét, nhấn mạnh: Ảnh của một hình là tập hợp toàn bộ các điểm ảnh tương ứng.
\end{itemize}

\subsubsection*{3. Hoạt động 3: Luyện tập}
\textbf{a) Mục tiêu:} Rèn luyện kỹ năng kiểm tra một quy tắc có phải là phép biến hình hay không; tìm tọa độ điểm ảnh trong mặt phẳng tọa độ $Oxy$.

\textbf{b) Nội dung:} Học sinh thực hiện Phiếu học tập số CĐ1 (Worksheet CĐ 1):
\begin{itemize}[leftmargin=0.6cm]
    \item \textbf{Bài 1:} Các quy tắc sau có phải là phép biến hình không? Giải thích rõ:
    \begin{enumerate}[leftmargin=0.5cm]
        \item Quy tắc $f_1$: Mỗi điểm $M(x; y)$ tương ứng với điểm $M'(2x; 3y - 1)$.
        \item Quy tắc $f_2$: Mỗi điểm $M(x; y)$ tương ứng với điểm $M'(x'; y')$ thỏa mãn $(x')^2 + (y')^2 = x^2 + y^2$.
    \end{enumerate}
    \item \textbf{Bài 2:} Trong mặt phẳng tọa độ $Oxy$, cho phép biến hình $f$ biến mỗi điểm $M(x; y)$ thành điểm $M'(x'; y')$ với công thức tọa độ:
    $$\begin{cases} x' = x + 1 \\ y' = 2y - 1 \end{cases}$$
    Tìm ảnh của các điểm $A(1; 2), B(-2; 0)$ và tìm điểm $C$ biết ảnh của nó là $C'(3; 5)$.
\end{itemize}

\textbf{c) Sản phẩm:} Lời giải chi tiết của học sinh:
\begin{itemize}[leftmargin=0.6cm]
    \item \textbf{Bài 1:}
    \begin{enumerate}[leftmargin=0.5cm]
        \item $f_1$ \textbf{là phép biến hình} vì với mỗi cặp $(x; y)$, các giá trị $x' = 2x$ và $y' = 3y - 1$ được xác định duy nhất.
        \item $f_2$ \textbf{không phải là phép biến hình} vì với một điểm $M(x; y) \neq (0; 0)$, có vô số điểm $M'$ nằm trên đường tròn tâm $O$ bán kính $\sqrt{x^2 + y^2}$, không thỏa mãn tính duy nhất.
    \end{enumerate}
    \item \textbf{Bài 2:}
    \begin{itemize}[leftmargin=0.5cm]
        \item Với $A(1; 2) \implies x' = 1 + 1 = 2; y' = 2 \cdot 2 - 1 = 3 \implies A'(2; 3)$.
        \item Với $B(-2; 0) \implies x' = -2 + 1 = -1; y' = 2 \cdot 0 - 1 = -1 \implies B'(-1; -1)$.
        \item Điểm $C(x; y)$ thỏa mãn: $\begin{cases} 3 = x + 1 \\ 5 = 2y - 1 \end{cases} \iff \begin{cases} x = 2 \\ 2y = 6 \end{cases} \iff \begin{cases} x = 2 \\ y = 3 \end{cases} \implies C(2; 3)$.
    \end{itemize}
\end{itemize}

\textbf{d) Tổ chức thực hiện:}
\begin{itemize}[leftmargin=0.6cm]
    \item \textit{Chuyển giao nhiệm vụ:} Học sinh làm bài độc lập trong 6 phút.
    \item \textit{Thực hiện nhiệm vụ (Hỗ trợ HS TB-Yếu):}
    \begin{scaffoldbox}{Giàn giáo sư phạm CĐ1.3: Hướng dẫn tìm điểm gốc khi biết điểm ảnh}
        Để tìm điểm gốc $C$ khi biết điểm ảnh $C'(x'; y')$, ta chỉ cần thay tọa độ của $C'$ vào vế trái của hệ công thức và giải phương trình bậc nhất một ẩn để tìm $x, y$.
    \end{scaffoldbox}
    \item \textit{Báo cáo, thảo luận:} 2 học sinh lên bảng trình bày, cả lớp đối chiếu và chấm chéo bài.
    \item \textit{Kết luận, nhận định:} Giáo viên đánh giá kỹ năng giải phương trình tọa độ của học sinh, khắc sâu quy tắc hai chiều: từ gốc tìm ảnh và từ ảnh truy ngược về gốc.
\end{itemize}

\subsubsection*{4. Hoạt động 4: Vận dụng}
\textbf{a) Mục tiêu:} Vận dụng khái niệm phép biến hình để hiểu nguyên lý xử lý hình ảnh trong Trí tuệ nhân tạo (AI) và đồ họa máy tính.

\textbf{b) Nội dung (Tích hợp Năng lực AI 11.C2.1):}
Giáo viên đặt bài toán: Khi một camera AI nhận diện khuôn mặt người dùng để mở khóa điện thoại:
\begin{itemize}[leftmargin=0.6cm]
    \item Khuôn mặt có thể bị nghiêng, xa hoặc gần camera. Làm thế nào để mô hình AI nhận ra đó vẫn là một khuôn mặt của bạn?
    \item AI đã sử dụng phép biến đổi tọa độ nào đối với các điểm mốc (landmarks) trên khuôn mặt (mắt, mũi, miệng)?
\end{itemize}

\textbf{c) Sản phẩm:} Lời giải và bài báo cáo của học sinh:
\begin{itemize}[leftmargin=0.6cm]
    \item Camera AI trích xuất các điểm mốc chính: 2 khóe mắt, chóp mũi, 2 khóe miệng (các điểm $M_1, M_2, \dots, M_n$).
    \item Thuật toán AI áp dụng các \textbf{phép biến hình tọa độ} (xoay góc, tịnh tiến dời gốc, co giãn tỉ lệ) để đưa tập hợp điểm ảnh về hệ quy chiếu chuẩn.
    \item So sánh khoảng cách tương đối giữa các điểm mốc để xác thực danh tính.
\end{itemize}

\textbf{d) Tổ chức thực hiện:}
\begin{itemize}[leftmargin=0.6cm]
    \item \textit{Chuyển giao nhiệm vụ:} Giáo viên trình chiếu slide hình ảnh nhận diện khuôn mặt qua FaceID/AI Camera.
    \item \textit{Thực hiện nhiệm vụ:} Học sinh thảo luận theo nhóm 4 em trong 3 phút.
    \item \textit{Báo cáo, thảo luận:} 1 nhóm đại diện phát biểu trình bày vai trò của phép biến hình trong xử lý ảnh AI.
    \item \textit{Kết luận, nhận định:} Giáo viên tổng kết Tiết CĐ1: Phép biến hình là nền tảng toán học của mọi phần mềm đồ họa và thuật toán thị giác máy tính AI hiện đại. Dặn dò chuẩn bị Tiết CĐ2 về Phép dời hình.
\end{itemize}

% =========================================================================
% TIẾT CĐ2
% =========================================================================
\vspace{10pt}
\subsection*{TIẾT CĐ2. PHÉP DỜI HÌNH VÀ CÁC TÍNH CHẤT CƠ BẢN}
\textit{(Tiết Chuyên đề 2 theo PPCT | Tuần 2)}

\subsubsection*{1. Hoạt động 1: Khởi động (Mở đầu)}
\textbf{a) Mục tiêu:} Nhận diện tính chất bảo toàn hình dạng và kích thước khi một vật thể chuyển động hoặc phản chiếu qua gương phẳng, dẫn dắt vào định nghĩa Phép dời hình.

\textbf{b) Nội dung:}
Giáo viên cho học sinh quan sát một miếng bìa hình tam giác $ABC$ cắt sẵn có $AB = 3\text{ cm}, AC = 4\text{ cm}, BC = 5\text{ cm}$.
\begin{itemize}[leftmargin=0.6cm]
    \item Khi trượt miếng bìa trên mặt bàn từ vị trí này sang vị trí khác, độ dài các cạnh của tam giác có bị thay đổi không?
    \item Nếu soi miếng bìa trước gương phẳng, tam giác ảnh $A'B'C'$ thu được có bằng tam giác ban đầu không? Khoảng cách giữa hai đỉnh bất kỳ có được giữ nguyên không?
\end{itemize}

\textbf{c) Sản phẩm:} Học sinh trả lời:
\begin{itemize}[leftmargin=0.6cm]
    \item Khi trượt miếng bìa, kích thước và hình dạng tam giác hoàn toàn không đổi: $A'B' = 3\text{ cm}, A'C' = 4\text{ cm}, B'C' = 5\text{ cm}$.
    \item Tam giác ảnh trong gương bằng tam giác gốc, khoảng cách giữa hai điểm bất kỳ luôn được giữ nguyên: $M'N' = MN$.
\end{itemize}

\textbf{d) Tổ chức thực hiện:}
\begin{itemize}[leftmargin=0.6cm]
    \item \textit{Chuyển giao nhiệm vụ:} Giáo viên thực hiện thao tác trực quan với miếng bìa tam giác vuông, gọi học sinh nhận xét.
    \item \textit{Thực hiện nhiệm vụ:} Học sinh quan sát, đối chiếu khoảng cách giữa các điểm.
    \item \textit{Báo cáo, thảo luận:} 2 học sinh nêu ý kiến về tính ''không làm biến dạng hình ảnh''.
    \item \textit{Kết luận, nhận định:} Giáo viên dẫn dắt: Phép biến hình mà không làm thay đổi khoảng cách giữa hai điểm bất kỳ được gọi là \textbf{Phép dời hình}.
\end{itemize}

\subsubsection*{2. Hoạt động 2: Hình thành kiến thức}

\paragraph{Nội dung 1: Khái niệm Phép dời hình}
\textbf{a) Mục tiêu:} Nhận biết và phát biểu chính xác định nghĩa phép dời hình; hiểu rõ điều kiện bảo toàn khoảng cách.

\textbf{b) Nội dung:}
\begin{itemize}[leftmargin=0.6cm]
    \item \textbf{Định nghĩa:} Phép dời hình là phép biến hình bảo toàn khoảng cách giữa hai điểm bất kỳ.
    \item Cụ thể: Nếu phép biến hình $f$ biến hai điểm $M, N$ bất kỳ lần lượt thành hai điểm $M', N'$ thì:
    $$M'N' = MN$$
    \item Phép đồng nhất $e$ là một phép dời hình (vì $M' \equiv M, N' \equiv N \implies M'N' = MN$).
\end{itemize}

\textbf{c) Sản phẩm:} Học sinh ghi nhớ định nghĩa toán học và công thức then chốt $M'N' = MN$.

\textbf{d) Tổ chức thực hiện:}
\begin{itemize}[leftmargin=0.6cm]
    \item \textit{Chuyển giao nhiệm vụ:} Giáo viên trình chiếu định nghĩa, yêu cầu học sinh viết công thức tổng quát vào vở.
    \item \textit{Thực hiện nhiệm vụ (Giàn giáo cho HS TB-Yếu):}
    \begin{scaffoldbox}{Giàn giáo sư phạm CĐ2.1: Ý nghĩa hình học của phép dời hình}
        \textbf{Từ khóa ghi nhớ:} \textbf{BẢO TOÀN KHOẢNG CÁCH}.
        Nghĩa là hình ảnh qua phép dời hình giống hệt hình ban đầu về kích thước và hình dạng (không bị phóng to, không bị thu nhỏ, không bị bóp méo).
    \end{scaffoldbox}
    \item \textit{Báo cáo, thảo luận:} 1 học sinh nhắc lại định nghĩa và nêu ví dụ trong đời sống (dịch chuyển đồ đạc, lật trang sách).
    \item \textit{Kết luận, nhận định:} Giáo viên chốt định nghĩa, chuyển sang nghiên cứu các tính chất hình học.
\end{itemize}

\paragraph{Nội dung 2: Các tính chất của Phép dời hình}
\textbf{a) Mục tiêu:} Nhận biết và vận dụng được 5 tính chất cơ bản của phép dời hình; hiểu bản chất biến một hình thành hình bằng nó.

\textbf{b) Nội dung:}
Phép dời hình có các tính chất cơ bản sau:
\begin{enumerate}[leftmargin=0.8cm]
    \item Biến ba điểm thẳng hàng thành ba điểm thẳng hàng và bảo toàn thứ tự giữa các điểm đó.
    \item Biến đường thẳng thành đường thẳng, biến tia thành tia, biến đoạn thẳng thành đoạn thẳng bằng nó.
    \item Biến tam giác thành tam giác bằng nó, biến góc thành góc bằng nó.
    \item Biến đường tròn bán kính $R$ thành đường tròn có cùng bán kính $R$.
    \item Biến đa giác $n$ cạnh thành đa giác $n$ cạnh có cùng diện tích và các cạnh tương ứng bằng nhau.
\end{enumerate}

\textbf{c) Sản phẩm:} Học sinh giải thích được:
\begin{itemize}[leftmargin=0.6cm]
    \item Nếu $B$ nằm giữa $A$ và $C$ thì $AB + BC = AC$. Do phép dời hình bảo toàn khoảng cách nên $A'B' + B'C' = A'C' \implies B'$ nằm giữa $A'$ và $C'$ (bảo toàn tính thẳng hàng và thứ tự).
    \item Tam giác $A'B'C'$ có 3 cạnh lần lượt bằng 3 cạnh tam giác $ABC \implies \triangle A'B'C' = \triangle ABC$ (theo trường hợp c-c-c).
    \item Đường tròn $(C)$ tâm $I$ bán kính $R$: Với mọi điểm $M \in (C) \iff IM = R \implies I'M' = R \iff M'$ thuộc đường tròn $(C')$ tâm $I'$ bán kính $R$.
\end{itemize}

\textbf{d) Tổ chức thực hiện:}
\begin{itemize}[leftmargin=0.6cm]
    \item \textit{Chuyển giao nhiệm vụ:} Giáo viên chia lớp thành 4 nhóm chuyên sâu, phân công mỗi nhóm thảo luận giải thích 1 tính chất (Nhóm 1: Tính chất ba điểm thẳng hàng; Nhóm 2: Đoạn thẳng; Nhóm 3: Tam giác; Nhóm 4: Đường tròn).
    \item \textit{Thực hiện nhiệm vụ:} Các nhóm thảo luận trong 4 phút, vẽ hình minh họa ra bảng phụ.
    \item \textit{Báo cáo, thảo luận:} Đại diện 4 nhóm trình bày bảng phụ. Các nhóm khác đặt câu hỏi chất vấn.
    \item \textit{Kết luận, nhận định:} Giáo viên chuẩn hóa 5 tính chất, nhấn mạnh: Hai hình được gọi là \textbf{bằng nhau} nếu có một phép dời hình biến hình này thành hình kia.
\end{itemize}

\subsubsection*{3. Hoạt động 3: Luyện tập}
\textbf{a) Mục tiêu:} Vận dụng định nghĩa và tính chất của phép dời hình để giải các bài toán chứng minh, tìm ảnh của đường thẳng và đường tròn.

\textbf{b) Nội dung:} Học sinh hoàn thành Phiếu học tập số CĐ2 (Worksheet CĐ 2):
\begin{itemize}[leftmargin=0.6cm]
    \item \textbf{Bài 1:} Cho tam giác $ABC$ vuông tại $A$ có $AB = 6\text{ cm}, AC = 8\text{ cm}$. Phép dời hình $f$ biến tam giác $ABC$ thành tam giác $A'B'C'$.
    \begin{enumerate}[leftmargin=0.5cm]
        \item Tính độ dài cạnh $B'C'$ và diện tích tam giác $A'B'C'$.
        \item Tam giác $A'B'C'$ có phải là tam giác vuông không? Vì sao?
    \end{enumerate}
    \item \textbf{Bài 2:} Trong mặt phẳng tọa độ $Oxy$, cho phép biến hình $f$ biến mỗi điểm $M(x; y)$ thành $M'(x'; y')$ có công thức:
    $$\begin{cases} x' = -x \\ y' = y \end{cases}$$
    \begin{enumerate}[leftmargin=0.5cm]
        \item Chứng minh $f$ là một phép dời hình.
        \item Tìm ảnh của đường tròn $(C): (x - 2)^2 + (y + 1)^2 = 9$ qua phép dời hình $f$.
    \end{enumerate}
\end{itemize}

\textbf{c) Sản phẩm:} Lời giải chi tiết của học sinh:
\begin{itemize}[leftmargin=0.6cm]
    \item \textbf{Bài 1:}
    \begin{enumerate}[leftmargin=0.5cm]
        \item Áp dụng định lý Pythagore trong $\triangle ABC$ vuông tại $A$:
        $$BC = \sqrt{AB^2 + AC^2} = \sqrt{6^2 + 8^2} = 10\text{ cm}$$
        Vì $f$ là phép dời hình nên $B'C' = BC = 10\text{ cm}$.\\
        Diện tích $\triangle A'B'C' = S_{\triangle ABC} = \dfrac{1}{2} \cdot AB \cdot AC = \dfrac{1}{2} \cdot 6 \cdot 8 = 24\text{ cm}^2$.
        \item Phép dời hình biến góc thành góc bằng nó nên $\widehat{B'A'C'} = \widehat{BAC} = 90^\circ \implies \triangle A'B'C'$ vuông tại $A'$.
    \end{enumerate}
    \item \textbf{Bài 2:}
    \begin{enumerate}[leftmargin=0.5cm]
        \item Lấy hai điểm bất kỳ $M(x_1; y_1)$ và $N(x_2; y_2)$. Khi đó qua $f$, ta có $M'(-x_1; y_1)$ và $N'(-x_2; y_2)$.\\
        Tính khoảng cách $MN$:
        $$MN = \sqrt{(x_2 - x_1)^2 + (y_2 - y_1)^2}$$
        Tính khoảng cách $M'N'$:
        $$M'N' = \sqrt{(-x_2 - (-x_1))^2 + (y_2 - y_1)^2} = \sqrt{(x_1 - x_2)^2 + (y_2 - y_1)^2} = MN$$
        Vì $M'N' = MN$ với mọi $M, N$ nên $f$ là một phép dời hình (đpcm).
        \item Đường tròn $(C)$ có tâm $I(2; -1)$ và bán kính $R = 3$.\\
        Qua phép dời hình $f$, tâm $I(2; -1)$ biến thành $I'(-2; -1)$ và bán kính không đổi $R' = R = 3$.\\
        Vậy phương trình đường tròn ảnh $(C')$ là:
        $$(x - (-2))^2 + (y - (-1))^2 = 3^2 \iff (x + 2)^2 + (y + 1)^2 = 9$$
    \end{enumerate}
\end{itemize}

\textbf{d) Tổ chức thực hiện:}
\begin{itemize}[leftmargin=0.6cm]
    \item \textit{Chuyển giao nhiệm vụ:} Học sinh làm việc cá nhân 7 phút, sau đó đổi bài chéo cặp đôi kiểm tra.
    \item \textit{Thực hiện nhiệm vụ (Giàn giáo sư phạm):}
    \begin{scaffoldbox}{Giàn giáo sư phạm CĐ2.2: Quy trình tìm ảnh của đường tròn qua phép dời hình}
        \textbf{Quy trình 3 bước chuẩn:}
        \begin{enumerate}[leftmargin=0.5cm]
            \item Xác định tâm $I(a; b)$ và bán kính $R$ của đường tròn gốc $(C)$.
            \item Dùng công thức phép biến hình để tìm tọa độ tâm ảnh $I'(a'; b')$.
            \item Giữ nguyên bán kính $R' = R$ và viết phương trình đường tròn $(C'): (x - a')^2 + (y - b')^2 = (R')^2$.
        \end{enumerate}
    \end{scaffoldbox}
    \item \textit{Báo cáo, thảo luận:} 2 học sinh lên bảng trình bày 2 bài. Lớp nhận xét chéo.
    \item \textit{Kết luận, nhận định:} Giáo viên nhận xét, nhấn mạnh: Trong phép dời hình, bán kính đường tròn tuyệt đối không thay đổi.
\end{itemize}

\subsubsection*{4. Hoạt động 4: Vận dụng (Giáo dục STEM \& Đồ họa văn hóa dân tộc)}
\textbf{a) Mục tiêu:} Vận dụng phép dời hình để giải thích nghệ thuật tạo hoa văn thổ cẩm của đồng bào dân tộc thiểu số A Lưới; phân tích vai trò kiểm soát của con người đối với AI tạo sinh đồ họa.

\textbf{b) Nội dung (Tích hợp STEM và Năng lực AI 11.A1.2):}
\begin{itemize}[leftmargin=0.6cm]
    \item \textbf{Nhiệm vụ 1 (STEM Văn hóa):} Quan sát dải hoa văn đường diềm trên trang phục thổ cẩm của đồng bào Pa Cô -- Tà Ôi (huyện A Lưới, Thừa Thiên Huế). Hãy chỉ ra một họa tiết cơ bản và các phép dời hình nào đã được người nghệ nhân dệt lặp lại để tạo thành dải hoa văn trang trí liên tục và cân đối?
    \item \textbf{Nhiệm vụ 2 (Năng lực AI 11.A1.2):} Hiện nay các công cụ AI tạo sinh hình ảnh (Midjourney, DALL-E) có thể tự động tạo ra hoa văn đối xứng trong vài giây. Tuy nhiên, tại sao con người vẫn cần kiểm tra và chỉnh sửa lại cấu trúc hình học của AI? Nếu AI bị ''ảo giác'' (sai lệch tỉ lệ khoảng cách), hậu quả trong in ấn dệt may sẽ như thế nào?
\end{itemize}

\textbf{c) Sản phẩm:} Báo cáo thu hoạch của học sinh:
\begin{itemize}[leftmargin=0.6cm]
    \item \textbf{Nhiệm vụ 1:} Họa tiết cơ bản là hình con thoi hoặc hoa văn mắt chim, hoa rừng. Nghệ nhân sử dụng phép tịnh tiến liên tiếp dọc theo mép vải kết hợp phép đối xứng qua trục để tạo ra sự đối xứng gương và lặp tuần hoàn, tạo cảm giác hài hòa, bền vững.
    \item \textbf{Nhiệm vụ 2:} AI tạo sinh dựa trên xác suất điểm ảnh nên thường xuyên làm lệch trục đối xứng, khoảng cách giữa các họa tiết không hoàn toàn bằng nhau ($M'N' \neq MN$). Nếu không có con người kiểm soát bằng toán học chính xác (các phép dời hình), mẫu in dệt công nghiệp sẽ bị méo mó, đường nối hoa văn bị đứt gãy, gây phế phẩm hàng loạt.
\end{itemize}

\textbf{d) Tổ chức thực hiện:}
\begin{itemize}[leftmargin=0.6cm]
    \item \textit{Chuyển giao nhiệm vụ:} Giáo viên trình chiếu mẫu hoa văn thổ cẩm A Lưới và mẫu hoa văn do AI tạo bị lỗi méo hình. Chia lớp thành 6 nhóm thảo luận trong 4 phút.
    \item \textit{Thực hiện nhiệm vụ:} Các nhóm phân tích hình ảnh, liên hệ giữa toán học và đời sống văn hóa.
    \item \textit{Báo cáo, thảo luận:} Đại diện 2 nhóm báo cáo, các nhóm khác phản biện.
    \item \textit{Kết luận, nhận định:} Giáo viên tổng kết toàn bộ bài học: Phép dời hình vừa là di sản toán học trong văn hóa truyền thống, vừa là thước đo chuẩn xác để con người làm chủ công nghệ trí tuệ nhân tạo. Giao bài tập về nhà trong Sách Chuyên đề học tập Toán 11.
\end{itemize}

% =========================================================================
% PHỤ LỤC
% =========================================================================
\newpage
\section*{IV. PHỤ LỤC HỌC LIỆU BỔ TRỢ}

\subsection*{1. Hệ thống Phiếu học tập (Worksheet) phân tầng bậc thang}

\begin{pedabox}{PHIẾU HỌC TẬP SỐ CĐ 1: PHÉP BIẾN HÌNH (Phân tầng 3 mức độ)}
\textbf{Họ và tên học sinh:} \dotfill \textbf{Lớp:} 11A\dots \ \textbf{Ngày:} \dotfill

\textbf{MỨC 1: NHẬN BIẾT (Điền khuyết có giàn giáo)}
\begin{enumerate}[leftmargin=0.5cm]
    \item Điền từ thích hợp vào chỗ trống:\\
    Phép biến hình trong mặt phẳng là quy tắc đặt tương ứng mỗi điểm $M$ với \dotfill điểm $M'$ của mặt phẳng đó. Điểm $M'$ được gọi là \dotfill của điểm $M$.
    \item Phép biến hình biến mỗi điểm $M$ thành chính nó được gọi là \dotfill.
\end{enumerate}

\textbf{MỨC 2: THÔNG HIỂU (Áp dụng công thức tọa độ)}
\begin{enumerate}[leftmargin=0.5cm]
    \item Trong mặt phẳng $Oxy$, cho phép biến hình $f$ biến điểm $M(x; y)$ thành $M'(x - 3; y + 4)$.\\
    a) Tìm ảnh của các điểm $O(0; 0)$ và $A(3; -4)$.\\
    b) Tìm tọa độ điểm gốc $B$ biết ảnh của nó là $B'(1; 2)$.\\
    \textit{Hướng dẫn giàn giáo:} Thay tọa độ $B'$ vào hệ: $x' = x - 3$ và $y' = y + 4$ để tìm $x, y$.
\end{enumerate}

\textbf{MỨC 3: VẬN DỤNG (Phân tích quy tắc)}
\begin{enumerate}[leftmargin=0.5cm]
    \item Cho trước đường thẳng $d$. Quy tắc biến mỗi điểm $M$ thành điểm $M'$ là hình chiếu vuông góc của $M$ trên $d$ có phải là phép biến hình không? Nếu lấy một đoạn thẳng $AB$ song song với $d$ thì ảnh của đoạn thẳng $AB$ là hình gì và có độ dài bằng bao nhiêu?
\end{enumerate}
\end{pedabox}

\vspace{8pt}

\begin{pedabox}{PHIẾU HỌC TẬP SỐ CĐ 2: PHÉP DỜI HÌNH \& TÍNH CHẤT (Phân tầng 3 mức độ)}
\textbf{MỨC 1: NHẬN BIẾT}
\begin{enumerate}[leftmargin=0.5cm]
    \item Phép dời hình là phép biến hình bảo toàn \dotfill giữa hai điểm bất kỳ.
    \item Điền ĐÚNG (Đ) hoặc SAI (S):\\
    a) Phép dời hình biến tam giác thành tam giác đồng dạng nhưng khác kích thước. (\dots)\\
    b) Phép dời hình biến đường tròn bán kính $R$ thành đường tròn có cùng bán kính $R$. (\dots)\\
    c) Phép dời hình bảo toàn tính thẳng hàng và thứ tự ba điểm. (\dots)
\end{enumerate}

\textbf{MỨC 2: THÔNG HIỂU}
\begin{enumerate}[leftmargin=0.5cm]
    \item Cho hình vuông $ABCD$ có cạnh bằng $4\text{ cm}$. Phép dời hình $f$ biến hình vuông $ABCD$ thành tứ giác $A'B'C'D'$. Tứ giác $A'B'C'D'$ là hình gì? Tính độ dài đường chéo $A'C'$ và diện tích tứ giác $A'B'C'D'$.
    \item Trong mặt phẳng tọa độ $Oxy$, cho phép biến hình $f$ biến mỗi điểm $M(x; y)$ thành $M'(y; x)$. Chứng minh $f$ là một phép dời hình.
\end{enumerate}

\textbf{MỨC 3: VẬN DỤNG}
\begin{enumerate}[leftmargin=0.5cm]
    \item Cho đường tròn $(C): x^2 + y^2 - 4x + 6y - 12 = 0$. Tìm phương trình đường tròn ảnh $(C')$ của $(C)$ qua phép dời hình biến mỗi điểm $M(x; y)$ thành $M'(x + 1; y - 2)$.\\
    \textit{Giàn giáo:} Bước 1: Đưa phương trình $(C)$ về dạng chính tắc $(x - a)^2 + (y - b)^2 = R^2$. Bước 2: Tìm tâm ảnh $I'$. Bước 3: Viết phương trình $(C')$.
\end{enumerate}
\end{pedabox}

\subsection*{2. Rubric đánh giá năng lực học sinh theo 4 mức độ}

\begin{table}[h!]
\centering
\small
\begin{tabularx}{\textwidth}{|p{2.8cm}|X|X|X|X|}
\hline
\textbf{Tiêu chí} & \textbf{Chưa đạt (Mức 1)} & \textbf{Đạt (Mức 2)} & \textbf{Khá (Mức 3)} & \textbf{Tốt (Mức 4)} \\
\hline
\textbf{1. Khái niệm Phép biến hình} & Chưa phân biệt được quy tắc có tính duy nhất hay không duy nhất. & Nhận biết được định nghĩa; tính được điểm ảnh của các điểm đặc biệt. & Xác định đúng ảnh của một hình cơ bản; tìm được điểm gốc khi biết ảnh. & Lập luận toán học chặt chẽ để chứng minh một quy tắc bất kỳ là phép biến hình. \\
\hline
\textbf{2. Khái niệm \& Tính chất Phép dời hình} & Nhầm lẫn phép dời hình với phép làm co giãn biến dạng hình ảnh. & Phát biểu được định nghĩa $M'N' = MN$; nhớ được các tính chất cơ bản. & Vận dụng tính chất tìm ảnh của đoạn thẳng, tam giác, đường tròn chuẩn xác. & Chứng minh được công thức khoảng cách bảo toàn trong hệ tọa độ $Oxy$; giải quyết bài toán cực trị hình học. \\
\hline
\textbf{3. Ứng dụng STEM \& Văn hóa} & Chưa nhận ra tính đối xứng và lặp lại của hoa văn thực tế. & Chỉ ra được hoa văn lặp lại trên trang phục thổ cẩm đơn giản. & Phân tích được phép biến hình tạo nên dải hoa văn đường diềm truyền thống. & Thiết kế được họa tiết hoa văn mới mang bản sắc dân tộc ứng dụng phép biến hình. \\
\hline
\textbf{4. Năng lực số \& Năng lực AI} & Chưa mở được GeoGebra; không biết công cụ dịch chuyển hình. & Thực hiện được thao tác vẽ và tịnh tiến điểm cơ bản trên phần mềm. & Khai thác thành thạo GeoGebra quan sát điểm ảnh; hiểu vai trò tọa độ trong AI. & Phân tích sâu sắc nguyên lý thị giác máy tính AI; chỉ ra tầm quan trọng của sự kiểm soát con người với AI. \\
\hline
\end{tabularx}
\end{table}

\end{document}